
\section{Introduction}
\label{sec:nonuniform-intro}

Auxiliary warping can generate stress fields that violate the sectional equilibrium equations \eqref{eq:sv-equilibrium}. 
This phenomenon was examined in detail by \cite{armero2021modeling}. 

Outline:
\begin{itemize}
    \item \cref{sec:nonuniform-prelim}: preliminaries and weighted inner products.
    \item \cref{sec:nonuniform-residual}: derivation of the residual.
    \item \cref{sec:nonuniform-enhanced}: construction of the minimal compatible correction space with one free parameter.
    \item \cref{sec:nonuniform-canonical}: canonical one-dimensional subspace selection.
    \item \cref{sec:nonuniform-solution}: condensed torsional response and general solution.
\end{itemize}


\section{Preliminaries}
\label{sec:nonuniform-prelim}

Consider moderate elastic deformations of an isotropic body with stresses of the form \eqref{eq:saint-venant-stress} so 
\begin{equation*}
    \bm{\sigma} = \mathbb{C}\bm{\varepsilon}
    \qquad\text{with}\qquad 
    \mathbb{C} = E \, \FrameBasis \otimes \FrameBasis + G \, \oneS
\end{equation*}
where \(E\) and \(G\) are Young's and the shear modulus, respectively, and the notation of \cref{sec:std-continuum} is employed. 
Define smooth, positive differential 2-forms \( \volG \) and \( \volE \)  on the cross-section \( \Section\) weighted by the shear \(G\) and Young's modulus \(E\), respectively:
\[
\begin{aligned}
    \volG &= G(\PlanePosition)\,\dA
    &\qquad
    \volE &= E(\PlanePosition)\,\dA \\
\end{aligned}
\]
where \(\dA\) is the Euclidean area form in the cross section plane \(\Section\). 
These serve as volume forms over the cross section and define divergence and  integration on the section $\Section$ with densities $\WeightE$ and $\WeightG$ relative to the standard area form (\cref{sec:integration}). 
Define the \(E\)-inner product, \(G\)-energy inner product, and \(G\)-Poisson bracket:
\begin{equation}
\left<u,\, v\right>_E \coloneq \int_{\Section} u \, v \, \volE,
\qquad
\left(u,\, v\right)_G \coloneq \int_{\Section} \nabla u \cdot \nabla v \, \volG,
\quad \text{and}\quad 
\left\{u,\, v\right\}_G \coloneq \int_{\Section} \nabla u \cdot \FrameBasis \times \nabla v\, \volG
\label{eq:def-products}
\end{equation}
of functions \(u,v\in L^{2}(\Section,\mathbb{R})\). 
The bracket \eqref{eq:def-products}$_3$ integrates the scalar triple product \([\nabla \phi,\, \FrameBasis,\, \PlanePosition] \coloneq \nabla \phi \cdot \FrameBasis \times \PlanePosition\)
\begin{equation*}
\tfrac{1}{2}\left\{\phi, r^2\right\}_G = \int [\nabla \phi,\, \FrameBasis,\, \PlanePosition]\, \mu_G
\end{equation*}


Consider the Type I warping beam theory where \(\WarpAxialShapes = \WarpTwistIesan\).
The Saint Venant function \(\WarpTwistIesan\) with regard to a point \(\bm{\pi} \in \Section\) is uniquely defined up to a constant as the Riesz vector 
\begin{equation}
    (\WarpTwistIesan,\, v)_G =  \frac{1}{2} \left\{ \| \PlanePosition - \bm{\pi}\|^2,\, v\right\}_G %= - [\nabla \psi,\, \FrameBasis, \PlanePosition]_G =
    \qquad \forall v \in \mathcal{F}(\Section)
    \label{eq:def-sv-riesz}
\end{equation}
or, equivalently (by \cref{prop:armero-81}), as the solution to the standard Neumann Laplace problem \eqref{eq:warp-sv-bvp}:
\begin{equation}
\left\{
\begin{array}{rll}
\DivG\left(\left.\nabla \WarpTwistIesan+\FrameBasis \times(\PlanePosition-\bm{\pi})\right.\right) &=0 & \text { in } \Section, \\
\nabla_{\bm{\nu}} \WarpTwistIesan &=-\FrameBasis\times(\PlanePosition-{\bm{\pi}}) \cdot \bm{\nu} & \text { on } \SectionBoundary, \\
\int_{\Section} \WarpTwistIesan \volE &= 0
\end{array}
\right.
% \label{eq:warp-sv-bvp}
% \quad\text{and}\quad
% \begin{array}{ll}
% \nabla \cdot\left(n_G \nabla \WarpFA \right)=n_E \WarpTwistIesan & \text { in } \Section, \\
% \nabla_{\bm{\nu}} \WarpFA =0 & \text { along } \partial \Section,
% \end{array}
\end{equation}
where \(\DivG\) is the divergence with respect to the volume form \(\mu_G\), as defined by \eqref{eq:def-div} (Box~\ref{box:recall}). 


\subsection{Pure Torsion}
Under pure torsion the rotation \(\bm{\Lambda}\)
can be represented with only a single parameter, say \(\vartheta\), so that:
\begin{equation}
\begin{aligned}
\bm{\Lambda} &= \operatorname{Exp} \vartheta \FrameBasis \\
&=\cos \vartheta \mathbf{1} + \sin \vartheta \, \FrameBasis^{\times} + (1-\cos \vartheta)  \FrameBasis \otimes \FrameBasis
\end{aligned}
\qquad\text{and}\qquad 
\begin{aligned}
\bm{\kappa}^{\times} &= \bm{\Lambda}^{\prime}\bm{\Lambda}^{\mathrm{t}} \\
&= \vartheta^{\prime} \, \FrameBasis^{\times}
\end{aligned}
\label{eq:rotation-field}
\end{equation}
from which \(\Twist = \vartheta'\). 
Note that under this motion the deformed (spatial) axial director \(\frameBasis\) remains equivalent to the reference axial director \(\FrameBasis\) (that is, \(\frameBasis = \FrameRotation \FrameBasis \equiv \FrameBasis\)). 
The compatible linearized strain tensor \(\hat{\StrainTensor}\) is characterized by the strain vector \(\hat{\bm{\varepsilon}}\) of \eqref{eq:linear-gl-strain-pi} which takes the form:
%
\begin{equation}
\begin{aligned}
\hat{\bm{\varepsilon}} 
&= \alpha' \WarpTwistIesan\,\FrameBasis + \Twist \FrameBasis \times (\PlanePosition-\tilde{\bm{\SectionLocation}}) + \alpha \nabla  \WarpTwistIesan \\
&= \alpha' \WarpTwistIesan\,\FrameBasis +  \Twist \, \bigl(\nabla \WarpTwistIesan + \FrameBasis\times(\PlanePosition - \tilde{\bm{\SectionLocation}})\bigr) - (\Twist - \alpha) \nabla \WarpTwistIesan
\end{aligned}
\label{eq:torsion-strain}
\end{equation}



\iftrue
% \section{Equilibrium Violations}
\section{Equilibrium residual}
\label{sec:nonuniform-residual}


% \begin{equation*}
% \FrameBasis \times \left(\SectionIntegral{\bar{\bm{r}}\otimes\bar{\bm{r}}} \right)^{-1}\SectionIntegral{\tilde{\varphi}\,\bar{\bm{r}}\times\FrameBasis}
% =-\FrameBasis^\times\tilde{\bm{\SectionLocation}}
% \end{equation*}

For any stress field satisfying \eqref{eq:saint-venant-stress} and the traction-free condition \(\StressTensor\,\bm{\nu}=\mathbf{0}\) on \(\BodyLateralSurface\),
the field equations reduce on each section to
\begin{subequations}
% \label{eq:sv-equilibrium}
\begin{align}
\bm{\tau}'&=\mathbf{0} && \text{in }\Section,\\
\DivS\bm{\tau}+\AxialStress'&=0 && \text{in }\Section, %\label{eq:sv-axial-equilibrium}
\\
\bm{\tau}\cdot\bm\nu&=0 && \text{on }\SectionBoundary.
\end{align}
\end{subequations}




Testing \eqref{eq:sv-axial-equilibrium} with any \(v\in H^1(\Section)\) and integrating by parts in \(\bm{r}\)
 gives the weak sectional equilibrium

\begin{equation}
\int_{\Section} G\,\PointShearStrain\cdot\nabla v\dA \;=\; -\int_{\Section} \partial_x(E\,\alpha'\WarpTwistIesan)\,v\dA
\qquad\forall v.
% \tag{5}
\label{eq:weak-nonuniform-residual}
\end{equation}

Since \(E\) and \(\WarpTwistIesan\) depend only on \(\bm{r}\),
\(\partial_x(E\alpha'\WarpTwistIesan)=E\,\alpha''\,\WarpTwistIesan\). 
Therefore the right
side is the functional

\begin{equation}
-\alpha''\int_{\Section} E\,\WarpTwistIesan\,v \dA.
\label{eq:nonuniform-functional-rhs}
\end{equation}

Therefore the equilibrium condition requires the mapping
\(v\mapsto\int_{\Section} G\,\PointShearStrain\cdot\nabla v\) to reproduce the mapping
\(v\mapsto-\alpha''\int_{\Section} E\WarpTwistIesan v\). 
%
Insert \eqref{eq:torsion-strain} into the left side of \eqref{eq:weak-nonuniform-residual} and apply \eqref{eq:def-sv-riesz}:

\begin{equation}
\int_{\Section} G\,\hat{\bm{\varepsilon}}_{\gamma}\cdot\nabla v
= (\alpha-\Twist)\int_{\Section} G\,\nabla\WarpTwistIesan\cdot\nabla v.
\label{eq:nonuniform-reduction}
\end{equation}

Therefore \eqref{eq:weak-nonuniform-residual} becomes

\begin{equation}
(\alpha-\Twist)\,L_1(v)+\alpha''\,L_2(v)=0 \qquad \forall v\in H^1(\Section)/\mathbb{R},
\label{eq:nonuniform-residual}
\end{equation}
where
% \begin{equation}
% (\alpha-\Twist)\underbrace{\int_{\Section} G\,\nabla\WarpTwistIesan\cdot\nabla v}_{=:L_1(v)}
% \;+\;\alpha''\underbrace{\int_{\Section} E\,\WarpTwistIesan\,v}_{=:L_2(v)}
% \;=\;0\qquad\forall v.
% \label{eq:nonuniform-residual}
% \end{equation}

% The residual \eqref{eq:nonuniform-residual} is
% a linear combination of two linear functionals on \(V\coloneq H^1(\Section)/\mathbb R\):

\begin{equation}
L_1(v) \coloneq (\WarpTwistIesan,\,v)_{G} %\int_{\Section} G\,\nabla\WarpTwistIesan\cdot\nabla v,
\qquad\text{and}\qquad
L_2(v) \coloneq \left<\WarpTwistIesan,\,v\right>_{E}.
%=\int_{\Section} E\,\WarpTwistIesan\,v.
\label{eq:nonuniform-functionals}
\end{equation}

\begin{Note}
TODO: Motivate modeling residual strain by a gradient field.
\end{Note}
To cancel \eqref{eq:nonuniform-residual} by adding a compatible correction shear field
\(\delta\tilde{\bm{\varepsilon}}_{\gamma}=\nabla\chi_c\), the potential \(\chi_c\) must satisfy:
%
\begin{equation}
\int_{\Section} G\,\nabla\chi_c\cdot\nabla v
= -(\alpha-\Twist) L_1(v)-\alpha''L_2(v)\quad\forall v.
\label{eq:nonuniform-functional}
\end{equation}
%
Define a function \(\WarpFA\)
\begin{equation}
\left\{
\begin{aligned}
\bigl( \DivG\nabla\WarpFA \bigr)\volG
  &=\WarpTwistIesan\,\volE.
&\quad\text{in }\Section, \\
\nabla_\nu \WarpFA &= 0 
&\quad \text{on } \SectionBoundary \\
\int_{\Section} \WarpFA \volE &= \AvgFA
\end{aligned}
\label{eq:warp-fa-bvp}
\right.
\end{equation}
considered by \cite{reissner1952nonuniform,armero2021modeling,genoese2013mixed}. 
Equivalently, \(\WarpFA\) is defined up to a constant \(\AvgFA\) as the Riesz representation  (\cref{prop:FA88})
\begin{equation}
(\WarpFA,\, v)_G = - \left<\WarpTwistIesan,\, v\right>_E
\qquad \forall v \in \mathcal{F}(\Section)
\label{eq:def-fa-riesz}
\end{equation}

Remark:

\begin{itemize}
\tightlist
\item
  \(L_1(\cdot)\) is represented by \(\WarpTwistIesan\):
  \(L_1(v)=\int_{\Section} G\nabla\WarpTwistIesan\cdot\nabla v\).
\item
  \(L_2(\cdot)\) is represented by \(\WarpFA\)
  through \eqref{eq:def-fa-riesz}.
\end{itemize}

Therefore the unique \(\chi_c\in V\) that satisfies \eqref{eq:nonuniform-functional} is

\begin{equation}
\chi_c = -(\alpha-\Twist)\WarpTwistIesan + \alpha''\WarpFA,
\label{eq:nonuniform-psi}
\end{equation}

and thus

\begin{equation}
\delta\tilde{\bm{\varepsilon}}_\gamma
= -(\alpha-\Twist)\nabla\WarpTwistIesan 
+ \alpha''\,\nabla\WarpFA.
\label{eq:equilibriating-strain}
\end{equation}

Therefore the residual is a linear combination of two independent functionals
\((L_1,L_2)\). The compatible representers of those functionals are
\((\nabla\WarpTwistIesan,\nabla\WarpFA)\), so the correcting field must be
their linear combination.


\section{Enhanced Subspace}
\label{sec:nonuniform-enhanced}

The minimal compatible correction space in view of \eqref{eq:nonuniform-residual} is

\[
\mathcal S \coloneq  \operatorname{span}\{\nabla\WarpTwistIesan,\ \nabla\WarpFA\}.
\]


The residual strain \eqref{eq:equilibriating-strain} can be injected by introducing an additional warping field with the shape \(\WarpFA\). 
However, in the absence of distributed warping loads \(\bm{f}_{\rm w}\),  \(\alpha'' \propto \alpha - \Twist\), in which case, the correction can be supplied within the section-local framework of \cref{sec:enhan-section}.
To this end, introduce the \(n_a=1\) enhanced shape with the form:
\begin{SaveEquation}{eq:nonuniform-shapes}
\EnhanShapes = \nabla\WarpTwistIesan + \EnhanConstFA\,\nabla\WarpFA
\end{SaveEquation}
% \[
% \boxed{\EnhanShapes_c(\bm r)\ \propto\ \nabla(\WarpTwistIesan+c\,\WarpFA),\qquad c\neq 0.}
% \]
where \(\EnhanConstFA \neq 0\) is required by the stability condition \eqref{eq:def-enhan-stability}. 
Define the enhanced potential

\[
\chi_c \coloneq  \WarpTwistIesan + \EnhanConstFA \,\WarpFA.
\]
Define constants:
%
\begin{subequations}\label{eq:J_defs}
\begin{alignat}{7}  % three column *pairs*: 7 = 2·3+1
\EJw        &\coloneq&\; &\langle\WarpTwistIesan,\WarpTwistIesan\rangle_E
           &\equiv&\;    &-(\WarpFA,\WarpTwistIesan)_G
           & &\;    & %[\nabla\WarpFA,\FrameBasis, \PlanePosition-\tilde{\bm{\SectionLocation}}]_G
           \label{eq:Jw}
           \\[4pt]
%
\GJv        &\coloneq&\;  &(\WarpTwistIesan,\WarpTwistIesan)_G
           & &\;           & %-[\nabla \WarpTwistIesan,\, \FrameBasis,\, \PlanePosition - \bm{\SectionLocation}]
           & &             &                              
           \label{eq:Jv}\\[4pt]
%
\Ja &\coloneq&\; &(\WarpFA,\WarpFA)_G
           &\equiv&\;    &-\langle\WarpTwistIesan,\WarpFA\rangle_E
           & &           &
\label{eq:Jsigma} \\[4pt]
%
\TwistRigidity &\coloneq&\; & \bigl(\WarpTwistIesan,\, \tfrac{1}{2} r^2\bigr)_G
           &+&    &\left(\tfrac{1}{2}r^2,\, \tfrac{1}{2}r^2\right)_G
           & &             &
\end{alignat}
\end{subequations}
% From your definitions

% \[
% (\WarpTwistIesan,\WarpTwistIesan)_G=\GJv,\qquad (\psi,\psi)_G=\Ja,\qquad (\psi,\WarpTwistIesan)_G=-\EJw,
% \]

% so the \(2\times2\) Gram data in \((\cdot,\cdot)_G\) is

% \[
% \begin{pmatrix}
% \GJv & -\EJw\\
% -\EJw & \Ja
% \end{pmatrix}.
% \]


The consistency condition \eqref{eq:def-enhan-residual} is

\begin{equation}
0=-(\Twist-\alpha)\,\tilde J(\EnhanConstFA) + \|\chi_c\|_G^2\,\EnhanParam,
\label{eq:nonuniform-consistency}
\end{equation}
%
where 
%
\begin{equation}
\tilde J(\EnhanConstFA)\coloneq (\WarpTwistIesan,\chi_c)_G
= (\WarpTwistIesan,\WarpTwistIesan)_G + \EnhanConstFA (\WarpTwistIesan,\WarpFA)_G
= \GJv - \EnhanConstFA\,\EJw,
\label{eq:def-J-enhan}
\end{equation}
%
and
%
\begin{equation}
\|\chi_c\|_G^2 \coloneq  (\chi_c,\chi_c)_G
= \GJv - 2\EnhanConstFA\,\EJw + \EnhanConstFA^2\Ja.
\label{eq:chi-norm}
\end{equation}
%
Hence
%
\begin{equation}
\EnhanParam=(\Twist-\alpha)\,\bar\eta(\EnhanConstFA),
\qquad 
\bar\eta(\EnhanConstFA)=\frac{\tilde J(\EnhanConstFA)}{\|\chi_c\|_G^2}
=\frac{\GJv-\EnhanConstFA\EJw}{\GJv-2\EnhanConstFA\EJw+\EnhanConstFA^2\Ja}.
\label{eq:nonuniform-eta}
\end{equation}




\paragraph{Summary.}

The correction space \(\mathcal S\) is reduced to one dimension by selecting a direction

\[
\nabla\chi_c \coloneq  \nabla(\WarpTwistIesan + \EnhanConstFA\,\WarpFA).
\]

Different values of \(\EnhanConstFA\) give different one-dimensional subspaces of \(\mathcal S\), and the reduced equilibrium determines only the amplitude of the chosen direction.
In this formulation, \eqref{eq:26} factors the correction as

\begin{equation*}
(\Twist-\alpha)\,\bar\eta\,\big(\nabla\WarpTwistIesan+\EnhanConstFA{}\nabla\WarpFA\big).
% \label{eq:nonuniform-correction-factor}
\end{equation*}

The correction vanishes when \(\Twist=\alpha\).


After reducing the two independent scalars \((\alpha-\Twist)\) and \(\alpha''\) to a single amplitude, the missing ratio required to represent the two-dimensional residual family \eqref{eq:nonuniform-residual} is absorbed into a scalar parameter:

\begin{equation}
\EnhanConstFA{}\ \text{is (up to sign) the effective ratio } \alpha''/(\Twist-\alpha)
\label{eq:enhan-const-ratio}
\end{equation}


\section{Cannonical Subspace}
\label{sec:nonuniform-canonical}

A clean normalization chooses \(\EnhanConstFA\) so that \(\chi_c\) is \(G\)-orthogonal
to \(\WarpFA\):

\[
(\chi_c,\WarpFA)_G=0
\quad\Longleftrightarrow\quad
(\WarpTwistIesan,\WarpFA)_G
+\EnhanConstFA (\WarpFA,\WarpFA)_G=0
\quad\Longleftrightarrow\quad
-\EJw+\EnhanConstFA \Ja=0,
\]

hence

\begin{equation}
\EnhanConstFA=\frac{\EJw}{\Ja}
= -\frac{(\WarpFA,\WarpTwistIesan)_G}{(\WarpFA,\WarpFA)_G}.
\label{eq:fa-lambda}
% \tag{9}
\end{equation}

With this choice,

\[
\|\chi_c\|_G^2=\tilde{J}(\EnhanConstFA)\quad\Longrightarrow\quad \bar\eta=1,
\]



\begin{remark}
This value is the unique \(\EnhanConstFA\) that satisfies either of the
equivalent characterizations:

\begin{enumerate}
\def\labelenumi{\arabic{enumi}.}
\tightlist
\item Gram--Schmidt \(G\)-orthogonalization:


\[
(\chi_c,\WarpFA)_G = -\EJw + \EnhanConstFA\Ja = 0.
\]

\item Minimum \(G\)-norm of the potential \(\chi_c\):

\[
\|\chi_c\|_G^2=\GJv-2\EnhanConstFA\EJw+\EnhanConstFA^2\Ja
\ \ \text{is minimized at}\ \ 
\frac{\dd}{\dd \EnhanConstFA}\|\chi_c\|_G^2 = -2\EJw+2\EnhanConstFA\Ja=0
\]

\end{enumerate}
\end{remark}


\section{General Solution}
\label{sec:nonuniform-solution}

Under pure constant torsion, the geometrically exact rod model (Box~\ref{box:exact-shear-svq}) furnishes the same response as the linearized model (Box~\ref{box:basic-shear}). A general solution for this case is derived in closed form.

\iftrue

\paragraph{Constitution.}
For \(\bm{\gamma} = \bm{0}\), \(\bm{n} = N \, \FrameBasis\). 
The moment has only a torsional component \(T=\bm{m}\cdot\FrameBasis\):
$$
\begin{aligned}
T 
&= \FrameBasis \cdot \int \PlanePosition \times \StressVector\, \dA \\
&= \GIo \Twist + \tfrac{1}{2}\{\WarpTwistIesan,\, r^2 \}_G \alpha 
 + \tfrac{1}{2} \left\{\WarpTwistIesan + \EnhanConstFA \WarpFA,\, r^2 \right\}_G \EnhanParam \\
&= \, \GIo\Twist - \GJv \alpha - (\GJv - \EnhanConstFA \, \EJw) \EnhanParam
\end{aligned}
\qquad\text{with}\qquad 
\begin{aligned}
\GIo &\triangleq \int_{\Section} r^2 \volG \\
\end{aligned}
$$
Using \eqref{eq:nonuniform-eta} furnishes:
\[
T = (\GIo - \bar{\eta}\tilde{J})\Twist - (\GJv - \bar{\eta}\tilde{J})\alpha
\]
The bi-moment is
\begin{equation}
B = \bm{w} = \int \WarpTwistIesan \FrameBasis \cdot \StressVector \dA =  \EJw \alpha'
\label{eq:nonuniform-bimoment}
\end{equation}
and the bi-shear is
\begin{equation*}
Q = \int \nabla \WarpTwistIesan \cdot \bm{\tau} \dA
= \tfrac{1}{2}\left\{\WarpTwistIesan,\, r^2\right\} \Twist 
+ \bigl(\WarpTwistIesan,\, \WarpTwistIesan \bigr)_G \alpha 
+ \bigl(\WarpTwistIesan,\, \WarpTwistIesan + \EnhanConstFA \WarpFA\bigr) \EnhanParam
= (-\GJv \Twist + \GJv \alpha + \tilde{J} \EnhanParam )
\end{equation*}


Substitute \(\EnhanParam\) from \eqref{eq:nonuniform-eta} and regroup:

\begin{equation}
Q
= \Big(\GJv(\alpha-\Twist)+\tilde{J}(\EnhanConstFA)(\Twist-\alpha)\bar\eta(\EnhanConstFA)\Big)
= (\alpha-\Twist)\Big(\GJv-\tilde{J}(\EnhanConstFA)\bar\eta(\EnhanConstFA)\Big).
\label{eq:nonuniform-bishear}
\end{equation}
\iffalse
Insert \eqref{eq:nonuniform-eta}:

\[
\GJv-\tilde J\,\bar\eta
=\GJv-\frac{\tilde{J}(\EnhanConstFA)^2}{\|\chi_c\|_G^2}.
% \tag{5}
\]

Using \eqref{eq:def-J-enhan}--\eqref{eq:chi-norm}, this simplifies to

\[
\GJv-\frac{\tilde{J}(\EnhanConstFA)^2}{\|\chi_c\|_G^2}
=\frac{\EnhanConstFA^2\big(\GJv\Ja-\EJw^2\big)}{\GJv-2\EnhanConstFA\,\EJw+\EnhanConstFA^2\Ja}.
% \tag{6}
\]

Therefore the condensed bi-shear is

\begin{equation}
Q= (\alpha-\Twist)\;
\frac{\EnhanConstFA^2\big(\GJv\Ja-(\EJw)^2\big)}{\GJv-2\EnhanConstFA\,\EJw+\EnhanConstFA^2\Ja}.
\label{eq:nonuniform-bishear}
%\tag{7}
\end{equation}

And bimoment equilibrium \(B'-Q=0\) \eqref{eq:finite-equilibrium}$_3$ becomes

\begin{equation}
\EJw\,\alpha''\;-\; (\alpha-\Twist)\;
\frac{\EnhanConstFA^2\big(\GJv\Ja-\EJw^2\big)}{\GJv-2\EnhanConstFA\,\EJw+\EnhanConstFA^2\Ja}=0.
\label{eq:bimoment-equilibrium}
% \tag{8}
\end{equation}


\fi 



%

\paragraph{Equilibrium.}
The moment equilibrium equation \eqref{eq:finite-equilibrium}$_2$ is stated in terms of the spatial moment \((\FrameRotation \bm{m})'\). 
However, for the rotation field \eqref{eq:rotation-field}, \(\FrameRotation\FrameBasis = \FrameBasis\) and the resulting equation exactly matches the linear theory:
\begin{equation}
\begin{array}{rlcl}
\FrameBasis \cdot (\FrameRotation \bm{m})' &= \bigl(\GIo \Twist - \GJv \alpha - (\GJv - \EnhanConstFA \, \EJw) \EnhanParam \bigr)' &=& 0 \\
\bm{w}' - \bm{v} &= \EJw \alpha'' - (\GJv (\alpha - \Twist) + \tilde{J} \EnhanParam) &=& 0
\end{array}
\label{eq:torsion-equilibrium}
\end{equation}
Note that \(\bm{x}' = \FrameBasis\) so that \(\FrameBasis\cdot \bm{x}' \times (\FrameRotation \bm{n})= \bm{0}\).
%
Substitution of \eqref{eq:fa-lambda} into \eqref{eq:enhan-strain} with \(\EnhanParam\) obtained in \eqref{eq:nonuniform-eta} furnishes
%
\begin{equation}
\begin{aligned}
\bm{\varepsilon} 
% &= \hat{\StrainVector} + \tilde{\StrainVector} \\
% &= \underbrace{\alpha' \WarpTwistIesan\, \FrameBasis + \Twist\, \FrameBasis \times \PlanePosition + \alpha \nabla \WarpTwistIesan}_{\hat{\StrainVector}} + \EnhanShape \EnhanParam  \\
&= \alpha' \WarpTwistIesan\, \FrameBasis + \Twist\, \FrameBasis \times \PlanePosition + \alpha \nabla \WarpTwistIesan + (\Twist - \alpha) \bar{\eta} (\nabla \WarpTwistIesan + \EnhanConstFA \nabla \WarpFA) \\
% &= \alpha' \WarpTwistIesan\, \FrameBasis + \Twist\, \FrameBasis \times \PlanePosition  - \alpha  \EnhanConstFA \nabla \WarpFA + \Twist (\nabla \WarpTwistIesan + \EnhanConstFA \nabla \WarpFA) \\
% &= \alpha' \WarpTwistIesan\, \FrameBasis + \Twist \, (\FrameBasis \times \PlanePosition  + \nabla \WarpTwistIesan) + (\Twist - \alpha)\EnhanConstFA \nabla \WarpFA \\
\end{aligned}
\label{eq:26}
\end{equation}
\begin{remark}
The shearing part of 
\eqref{eq:torsion-strain} agrees with the strain associated with the RBV formulation \citep[eq. 108]{armero2021modeling} while \eqref{eq:26} is that of the Mixed formulation \citep[eq. 161]{armero2021modeling}.
\end{remark}
%
\paragraph{Solution.}
Integrating the moment equilibrium equation \eqref{eq:torsion-equilibrium}$_1$ furnishes:
\begin{equation*}
    \alpha (\reflenx)= \hat{\eta }\bigl(
     \vartheta^{\prime}(\reflenx) - C\bigr)
     \quad\text{with}\quad 
     \hat{\eta} = \frac{\GIo - \bar{\eta}\tilde{J}}{\GJv - \tilde{J}}
\end{equation*}
where \(C\) is a constant of integration. 
This produces the bimoment \eqref{eq:nonuniform-bimoment} and bishear \eqref{eq:nonuniform-bishear} expressed in terms of the twist rate \(\Twist\):
\begin{equation*}
B = \EJw \hat{\eta} \Twist'(\reflenx) 
\qquad\text{and}\qquad
Q = (\GJv + \bar{\eta}\tilde{J})(1 - \hat{\eta}) \Twist(\reflenx)
\end{equation*}
and equilibrium of the bimoment \eqref{eq:torsion-equilibrium}$_2$ becomes:
\begin{equation}
\EJw\, \hat{\eta}\, \Twist'' - (\GJv + \tilde{J}\bar{\eta})(1 - \hat{\eta}) \Twist = 0
\end{equation}
%
In standard form:
\begin{equation*}
\vartheta''' - \lambda^2 \vartheta' = 0,
\qquad\text{with}\quad
\lambda \triangleq \sqrt{\frac{\GJv - \bar{\eta}\tilde{J}}{\EJw}\frac{1 - \hat{\eta}}{\hat{\eta}}}
\end{equation*}
%
This furnishes the general solution with four free constants \(\TwistRxn\) and \(C_k\):
\begin{equation}
\begin{aligned}
\vartheta(\reflenx)= \frac{\TwistRxn}{\lambda \hat{\eta} \, \TwistRigidity}\left(
 C_1 + \lambda\hat{\eta} \, \reflenx +  C_2 \cosh \lambda \reflenx + C_3 \sinh \lambda \reflenx
\right)
\end{aligned}
\label{eq:0028-gen-rot}
\end{equation}

with boundary conditions:

% The beam ends can have either an imposed rotation $\vartheta=\vartheta_0$ or an imposed torque $T=\TwistRxn$. 
% At the same time they have either an imposed warping $\vartheta'=\vartheta_0^{\prime}$ or an imposed bimoment $B=B_0$. 
\begin{enumerate}
\item $\mathrm{D}_{\vartheta}$ twisting angle restrained $\vartheta=0$
\item $\mathrm{D}_{\alpha}$ warping restrained $\alpha=0$ or $\vartheta^{\prime}=0$
\item Warping free $B \,(= \alpha') =0$ or $\theta^{\prime \prime}=0$
\item Twisting free $\vartheta^{\prime \prime \prime}-\lambda^2 \vartheta^{\prime}=0$ %$\qquad\implies T=B'$
\end{enumerate}
\fi

\paragraph{Parameter Constraints}


\begin{itemize}
\tightlist
\item
  Well-posedness of the condensation requires
  \(\|\chi_c\|_G^2>0\). Since


\[\|\chi_c\|_G^2
= \Ja\Big(\EnhanConstFA-\frac{\EJw}{\Ja}\Big)^2+\Big(\GJv-\frac{\EJw^2}{\Ja}\Big),\]

this is \(>0\) for all \(\EnhanConstFA\) provided the Gram determinant
\(\GJv\Ja-\EJw^2>0\) (i.e.~\(\WarpTwistIesan\) and \(\WarpFA\) are
\(G\)-independent). 
For the degenerate case \(\GJv\Ja=\EJw^2\), the energy \(\|\chi_c\|_G^2\)
vanishes at 
\[
\EnhanConstFA=\EJw/\Ja
\]

\item Non-degeneracy of the resulting beam coupling:
after the condensation, the bi-shear reduces to a
coefficient times \((\alpha-\Twist)\). 
One finds 

\[
Q = (\alpha-\Twist)\Big(\GJv-\frac{\tilde J(\EnhanConstFA)^2}{\|\chi_c\|_G^2}\Big)
   = (\alpha-\Twist)\,\frac{\EnhanConstFA^2(\GJv\Ja-\EJw^2)}{\GJv-2\EnhanConstFA\EJw+\EnhanConstFA^2\Ja}.
\]

In particular, \(\EnhanConstFA=0\Rightarrow Q\equiv 0\), i.e.~this eliminates the
bimoment coupling. 
Therefore \(\EnhanConstFA=0\) is mathematically admissible but
physically degenerates the Vlasov effect, recovering the TWKV formulation.
\end{itemize}
